% ===== PRÉAMBULE CANONIQUE — à coller VERBATIM en tête de CHAQUE .tex =====
\documentclass[11pt,a4paper]{article}
\usepackage[utf8]{inputenc}\usepackage[T1]{fontenc}\usepackage{lmodern}
\usepackage[french]{babel}
\usepackage{amsmath,amssymb,mathtools}\usepackage{icomma}
\usepackage[version=4]{mhchem}          % \ce{...} : équations & formules
\usepackage{siunitx}\sisetup{locale=FR} % \qty{}{}, \si{}, \ang{}
\usepackage{chemfig}                    % \chemfig{...} : structures développées
\usepackage{xcolor}\usepackage{colortbl}
\usepackage{tikz}\usepackage{pgfplots}\pgfplotsset{compat=1.18}
\usepackage[most]{tcolorbox}\usepackage{enumitem}\usepackage{array,booktabs}
\usepackage[a4paper,margin=2.2cm]{geometry}
\usetikzlibrary{arrows.meta,calc,angles,quotes,positioning,patterns,backgrounds,shapes.geometric}
\definecolor{primary}{RGB}{222,49,99}\definecolor{primarydark}{RGB}{150,22,55}
\definecolor{defcol}{RGB}{0,114,178}\definecolor{methcol}{RGB}{52,92,148}
\definecolor{excol}{RGB}{0,135,90}\definecolor{attncol}{RGB}{200,40,55}
\tcbset{boxbase/.style={enhanced,breakable,boxrule=0.9pt,arc=2.5pt,
  left=3mm,right=3mm,top=2.3mm,bottom=2.3mm,fonttitle=\bfseries,coltitle=white}}
% --- boîtes TUTORIEL (noms EXACTS, ne pas renommer) ---
\newtcolorbox{cle}[1][]{boxbase,colback=defcol!6,colframe=defcol,
  title={\textsc{À retenir}\ifx\relax#1\relax\relax\else\textmd{ : \itshape #1}\fi}}
\newtcolorbox{methode}[1][]{boxbase,colback=methcol!7,colframe=methcol,sharp corners,
  title={\textsc{Méthode}\ifx\relax#1\relax\relax\else\textmd{ --- \itshape #1}\fi}}
\newtcolorbox{exemplebox}[1][]{boxbase,colback=excol!5,colframe=excol,
  title={\textsc{Exemple}\ifx\relax#1\relax\relax\else\textmd{ : \itshape #1}\fi}}
\newtcolorbox{piege}[1][]{boxbase,colback=attncol!6,colframe=attncol,
  title={\textsc{Piège}\ifx\relax#1\relax\relax\else\textmd{ : \itshape #1}\fi}}
\newtcolorbox{remarque}[1][]{boxbase,colback=black!4,colframe=black!45,
  title={\textsc{Remarque}\ifx\relax#1\relax\relax\else\textmd{ : \itshape #1}\fi}}
% --- boîtes EXERCICE / CORRIGÉ (noms EXACTS, auto-comptées) ---
\newtcolorbox[auto counter]{exo}[1][]{boxbase,colback=primary!4,colframe=primary,
  title={\textsc{Exercice}~\thetcbcounter\ifx\relax#1\relax\relax\else\textmd{ --- \itshape #1}\fi}}
\newtcolorbox[auto counter]{corrige}[1][]{boxbase,colback=excol!4,colframe=excol,
  title={\textsc{Corrigé}~\thetcbcounter\ifx\relax#1\relax\relax\else\textmd{ --- \itshape #1}\fi}}
\newcommand{\R}{\mathbb{R}}\newcommand{\N}{\mathbb{N}}\newcommand{\Z}{\mathbb{Z}}\newcommand{\ds}{\displaystyle}
\begin{document}
\begin{corrige}[Classer les carbones des trois isomères de $\ce{C5H12}$]
\begin{enumerate}[label=\alph*)]
  \item \textbf{pentane} \ce{CH3-CH2-CH2-CH2-CH3} : les deux bouts sont primaires, les trois carbones internes secondaires $\Rightarrow$ \textbf{2 primaires, 3 secondaires} (0 tertiaire, 0 quaternaire).
  \item \textbf{2-méthylbutane} : $\Rightarrow$ \textbf{3 primaires, 1 secondaire, 1 tertiaire} (0 quaternaire).
  \item \textbf{néopentane} : le carbone central est lié à $4$ méthyles $\Rightarrow$ \textbf{4 primaires, 1 quaternaire} (0 secondaire, 0 tertiaire).
\end{enumerate}
Remarque : plus la molécule est ramifiée, plus elle contient de carbones de classe élevée.
\end{corrige}

\begin{corrige}[Lire les carbones d'une molécule polyfonctionnelle]
Numérotons \ce{C1H3-C2O-C3H2-C4H=C5H-C6HO}.
\begin{enumerate}[label=\alph*)]
  \item $6$~C ; H $=3+0+2+1+1+1=8$ ; O $=2$ $\Rightarrow \ce{C6H8O2}$.
  \item Carbones à $2$~H : \textbf{un seul} ($\ce{C3}$, le $\ce{-CH2-}$). Carbones à $1$~H : \textbf{trois} ($\ce{C4}$, $\ce{C5}$ de la double liaison, et $\ce{C6}$ de l'aldéhyde). Carbones sans H : \textbf{un} ($\ce{C2}$). Le $\ce{C1}$ ($\ce{CH3}$) porte $3$~H.
  \item C'est $\ce{C2}$ : c'est le carbone d'une \textbf{cétone}, lié à $\ce{C1}$ (1), à $\ce{C3}$ (1) et doublement à O (2), soit $4$ liaisons vers des atomes autres que H $\Rightarrow$ $n_{\text{H}}=4-4=0$.
\end{enumerate}
\end{corrige}

\begin{corrige}[Identifier un isomère par la description de ses carbones]
\begin{enumerate}[label=\alph*)]
  \item Quatre carbones liés à $3$~H sont quatre groupements $\ce{-CH3}$. Pour $\ce{C5}$, il reste un cinquième carbone au centre : c'est le \textbf{2,2-diméthylpropane (néopentane)}.
  \item Les quatre $\ce{CH3}$ occupent déjà les quatre liaisons du carbone central : il ne lui reste \textbf{aucune} liaison pour un H $\Rightarrow$ $0$~H, c'est un carbone \textbf{quaternaire}.
  \item Décompte : \textbf{4 carbones primaires} (les $\ce{CH3}$) et \textbf{1 carbone quaternaire}.
\end{enumerate}
\end{corrige}

\begin{corrige}[Tous les alcènes de $\ce{C4H8}$]
\begin{enumerate}[label=\alph*)]
  \item $I=\frac{2\cdot4+2-8}{2}=\frac{2}{2}=1$ : une insaturation, compatible avec une double liaison $\ce{C=C}$.
  \item Les trois alcènes de constitution : \textbf{but-1-ène} \ce{CH2=CH-CH2-CH3}, \textbf{but-2-ène} \ce{CH3-CH=CH-CH3} et \textbf{2-méthylprop-1-ène} \ce{CH2=C(CH3)-CH3} (isobutène).
  \item Seul le \textbf{but-2-ène} présente l'isomérie cis/trans (chaque carbone de la $\ce{C=C}$ porte un $\ce{CH3}$ et un $\ce{H}$). On obtient donc \textbf{4} structures d'alcènes distinctes : but-1-ène, \emph{cis}-but-2-ène, \emph{trans}-but-2-ène et 2-méthylprop-1-ène.
\end{enumerate}
\end{corrige}

\end{document}
